\subsection{Variables cualitativas}
\label{sec:eda-cualitativas}

Las categóricas se reportan en dos bloques por unidad de análisis:
\texttt{periodo} y \texttt{llovio} solo existen a resolución horaria; las
indicadoras de excedencia derivan del MDA8 y solo existen a resolución
diaria. Mezclarlas replicaría cada valor diario 24 veces. Las bases son
$n = 640{,}583$ registros hora-estación y $n = 26{,}691$ día-estación.

\subsubsection{Distribución de frecuencia}

\Cref{tab:freq-horaria} reúne las variables horarias. Con 2025 completo las
cuatro temporadas quedan equilibradas (24.79\,\% a 25.16\,\%), de modo que
ningún contraste estacional es artefacto de la ventana de observación.
\texttt{periodo} es uniforme por construcción: cuatro bloques de 6 horas
sobre una rejilla horaria completa. \texttt{llovio} es la única horaria con
faltantes (19,360 registros, 3.02\,\%), que se conservan como categoría
propia en vez de asignarse a ``no llovió''; llueve en 1.49\,\% de las horas
con medición. \texttt{tipo\_dia} sustituye a la binaria de fin de semana
porque separa sábado de domingo, cuyos patrones vehiculares difieren, y
aísla los 71 días festivos del periodo.

\begin{table}[t]
  \centering
  \caption{Frecuencia absoluta y relativa de las variables cualitativas a
    resolución hora-estación ($n = 640{,}583$). Los porcentajes son sobre
    registros con dato válido.}
  \label{tab:freq-horaria}
  \small
  \begin{tabular}{llrr}
    \toprule
    Variable & Nivel & Frecuencia & \% \\
    \midrule
    \multirow{4}{*}{\texttt{temporada}}
                                  & Invierno  &   158,783 &  24.79 \\
                                  & Primavera &   161,184 &  25.16 \\
                                  & Verano    &   161,184 &  25.16 \\
                                  & Otoño     &   159,432 &  24.89 \\
    \midrule
    \multirow{4}{*}{\texttt{periodo}}
                                  & Madrugada &   160,145 &  25.00 \\
                                  & Mañana    &   160,146 &  25.00 \\
                                  & Tarde     &   160,146 &  25.00 \\
                                  & Noche     &   160,146 &  25.00 \\
    \midrule
    \multirow{4}{*}{\texttt{tipo\_dia}}
                                  & Laborable &   434,879 &  67.89 \\
                                  & Sábado    &    88,392 &  13.80 \\
                                  & Domingo   &    88,728 &  13.85 \\
                                  & Festivo   &    28,584 &   4.46 \\
    \midrule
    \multirow{3}{*}{\texttt{llovio}}
                                  & No (0)    &   611,952 &  98.51 \\
                                  & Sí (1)    &     9,271 &   1.49 \\
                                  & Sin dato  &    19,360 &    --- \\
    \bottomrule
  \end{tabular}
\end{table}

\Cref{tab:freq-diaria} presenta el bloque diario. Las tres indicadoras de
MDA8 comparten los mismos 1,591 registros sin dato (5.96\,\%), los
días-estación cuyo MDA8 no alcanzó la completitud de 18 de 24 ventanas;
\texttt{excede\_1h} no los comparte porque se calcula sobre el máximo
horario.

\begin{table}[t]
  \centering
  \caption{Frecuencia absoluta y relativa de las indicadoras de excedencia a
    resolución día-estación ($n = 26{,}691$; válidos $= 25{,}100$ para las
    indicadoras de MDA8).}
  \label{tab:freq-diaria}
  \small
  \begin{tabular}{llrr}
    \toprule
    Variable & Nivel & Frecuencia & \% \\
    \midrule
    \multirow{3}{*}{\texttt{excede\_anio\_1} (65 ppb)}
                                        & No excede &  22,275 &  88.75 \\
                                        & Excede    &   2,825 &  11.25 \\
                                        & Sin dato  &   1,591 &    --- \\
    \midrule
    \multirow{3}{*}{\texttt{excede\_anio\_3} (60 ppb)}
                                        & No excede &  20,864 &  83.12 \\
                                        & Excede    &   4,236 &  16.88 \\
                                        & Sin dato  &   1,591 &    --- \\
    \midrule
    \multirow{3}{*}{\texttt{excede\_anio\_5} (51 ppb)}
                                        & No excede &  17,079 &  68.04 \\
                                        & Excede    &   8,021 &  31.96 \\
                                        & Sin dato  &   1,591 &    --- \\
    \midrule
    \multirow{2}{*}{\texttt{excede\_1h} (90 ppb)}
                                        & No excede &  24,497 &  91.78 \\
                                        & Excede    &   2,194 &   8.22 \\
    \bottomrule
  \end{tabular}
\end{table}

\subsubsection{La estación de monitoreo}

Catorce de las 15 estaciones registran 43,824 observaciones cada una, la
rejilla horaria completa de 2021-01-01 a 2025-12-31; \texttt{NO3} registra
27,047 porque inicia operación el 1 de diciembre de 2022. La diferencia es de
cobertura, no de calidad: dentro de su ventana \texttt{NO3} también tiene
rejilla completa y su porcentaje de horas con \ce{O3} medido está dentro del
rango del resto de la red. \texttt{periodo}, la única
ordinal, tiene mediana \texttt{tarde} sobre las 611,601 horas con \ce{O3}
medido, porque los faltantes se concentran de noche.

\subsubsection{Balance de clases}

\Cref{fig:eda-cual-excedencias} visualiza el balance de las tres indicadoras
de MDA8. Al umbral vigente de 65 ppb las excedencias son 11.25\,\% de los
días-estación válidos: un clasificador que prediga siempre la clase
mayoritaria alcanza 88.75\,\% de aciertos sin capturar el fenómeno, y la
exactitud queda descartada como métrica en favor de sensibilidad,
especificidad y área bajo la curva ROC. El umbral de 51 ppb, con
31.96\,\%, ofrece un balance sustancialmente mejor y es el escenario
principal de modelado.

\begin{figure}[t]
  \centering
  \includegraphics[width=\linewidth]{eda_cual_excedencias.pdf}
  \caption{Balance de clases de las indicadoras de excedencia del MDA8 para los
    tres umbrales de cumplimiento gradual de la NOM-020-SSA1-2021
    (\cref{tab:nom-limites}).}
  \label{fig:eda-cual-excedencias}
\end{figure}
